\section{Frequency-Report Scoring Rules}

This section applies the setup to five count-loss rules.
The first two rules have simple analytic identification bounds and are the primary theoretical objects in the paper.
The remaining rules have exact inverse-region representations that are useful for computational comparison, but their sharp bounds are obtained from those representations rather than from comparably simple main-text formulas.

\subsection{Primary Rules with Analytic Bounds}

\subsubsection{Quadratic-Distance Frequency Scoring}

The main analytical rule is quadratic-distance frequency scoring.
It pays less when the reported count vector is far from the realized count vector in squared Euclidean distance:
\[
S_Q(r,\omega)=a-b\|r-\omega\|_2^2,
\quad b>0.
\]
This rule is connected in spirit to quadratic scoring rules, but it is not a standard probability score.
It is a frequency-report count-loss rule: subjects report counts, and the penalty is the squared difference between reported and realized counts.
Write \(P_Q\) for the belief region induced by this rule.

Given beliefs \(p\), the expected count vector is \(np\).
The key forward fact is that quadratic-distance scoring asks the subject to report the feasible integer count vector closest to \(np\).
The inverse question is therefore geometric: for which beliefs is the observed report \(r\) the closest feasible count vector to \(np\)?

\begin{proposition}[Quadratic distance: reports, regions, and bounds]
For every \(p\in\DeltaK\),
\[
R_{S_Q}(p)
=
\argmin_{s\in\OmegaN}\|s-np\|_2^2.
\]
For an observed report \(r\in\OmegaN\), the belief identification region is
\[
P_Q(r)
=
\left\{
p\in\DeltaK:
n(p_i-p_j)\le r_i-r_j+1
\quad
\forall i,j\text{ with }r_j>0
\right\}.
\]
Let
\[
m(r)=|\{j:r_j>0\}|
\]
be the number of categories with positive reported counts.
The sharp coordinate intervals are
\[
P_{Q,i}(r)=
\begin{cases}
\displaystyle
\left[0,\frac{m(r)}{n(m(r)+1)}\right],
& r_i=0,\\[1.2em]
\displaystyle
\left[
\frac{r_i-1}{n}+\frac{1}{nk},
\frac{r_i+1}{n}-\frac{1}{n\,m(r)}
\right],
& r_i>0.
\end{cases}
\]
Projection ties are included.
\end{proposition}

The inequalities compare the observed report \(r\) with reports obtained by moving one count from a positive reported category \(j\) to another category \(i\).
If all such one-count moves weakly increase squared distance from \(np\), then \(r\) is a closest feasible report.
Thus \(P_Q(r)\) is a polytope in the belief simplex.
The proof is in Appendix \ref{app:squared-proof}.

\paragraph{Probability and mean bounds.}
The coordinate bounds are closed form.
For numerical outcomes \(x_1,\ldots,x_k\), sharp mean bounds are obtained by optimizing over the same polytope:
\[
\min_{p\in P_Q(r)}\sum_{i=1}^k p_i x_i
\quad\text{and}\quad
\max_{p\in P_Q(r)}\sum_{i=1}^k p_i x_i.
\]
These are linear programs.
For coordinate probabilities, the displayed intervals give the corresponding linear-program solutions explicitly.
This is the main practical advantage of quadratic-distance frequency scoring: the same transparent region gives both probability bounds and bounds on expected values.

\paragraph{Practical interpretation.}
Quadratic distance is mean-oriented: it links a report \(r\) to beliefs for which \(r\) is the integer projection of \(np\).
Compared with exact matching, the payment varies smoothly with distance rather than only with exact equality.
The main implementation cost is that direct monetary quadratic-distance payments generally require risk neutrality unless the score is implemented probabilistically, as discussed below.

\subsubsection{Discrete-Metric Frequency Scoring}

The second primary rule is the frequency-guessing mechanism of Schlag and Tremewan \cite{SchlagTremewanSimple}.
It is the count-vector discrete metric:
\[
D_0(r,\omega)=\1\{r\ne\omega\}.
\]
Equivalently, a fixed prize is paid if and only if the realized count vector is exactly equal to the report:
\[
q_0(r,\omega)=\1\{r=\omega\}.
\]
Under risk neutrality, this is equivalent to the affine score
\[
S_0(r,\omega)=a+b\1\{r=\omega\},
\quad b>0.
\]
The fixed-prize formulation is useful for implementation and risk aversion; the affine-score formulation is convenient for comparison with the distance rules.
Write \(R_0\) and \(P_0\) for the report correspondence and belief region induced by this rule.

The mechanism is simple: a subject with beliefs \(p\) maximizes the probability of an exact match,
\[
\Pr_p(\omega=r)
=
\frac{n!}{r_1!\cdots r_k!}\prod_{i=1}^k p_i^{r_i}.
\]
Thus the forward object is the multinomial mode set.
The inverse object is the set of beliefs for which the observed report is one of those modes.

\begin{proposition}[Discrete metric: reports, regions, and bounds]
Let \(r\in\OmegaN\), and allow ties in the multinomial mode.
Then \(r\) is optimal under discrete-metric frequency scoring if and only if \(r\) is a multinomial mode:
\[
R_0(p)
=
\argmax_{s\in\OmegaN}\Pr_p(\omega=s).
\]
The belief identification region is
\[
P_0(r)
=
\left\{
p\in\DeltaK:
r_jp_i\le (r_i+1)p_j
\quad
\forall i,j\text{ with }r_j>0
\right\}.
\]
For each coordinate, the sharp identification interval is
\[
P_{0,i}(r)
=
\left[
\frac{r_i}{n+k-1},
\frac{r_i+1}{n+1}
\right].
\]
\end{proposition}

The inequalities say that no one-count transfer from a reported category \(j\) to another category \(i\) can increase the multinomial probability of the report.
Weak inequalities include ties.
The formula also handles boundaries: if \(r_i=0\), the lower bound is \(0\); if \(r_i=n\), the upper bound is \(1\).
These are the frequency-guessing bounds in Schlag and Tremewan \cite[Proposition 1]{SchlagTremewanSimple}; they are restated here to make the informational-efficiency comparison self-contained.
The proof is in Appendix \ref{app:exact-match-proof}.

\paragraph{Probability and mean bounds.}
The discrete-metric rule gives closed-form coordinate intervals.
For linear functionals such as means, sharp bounds use the full joint region, not just the coordinate intervals:
\[
\min_{p\in P_0(r)}\sum_{i=1}^k p_i x_i
\quad\text{and}\quad
\max_{p\in P_0(r)}\sum_{i=1}^k p_i x_i.
\]
These are linear programs because \(P_0(r)\) is defined by linear inequalities.
Thus discrete-metric frequency scoring is analytically convenient for individual probability bounds, while mean inference still requires using the joint mode region.

\paragraph{Practical interpretation.}
Discrete-metric frequency scoring is transparent and robust to risk aversion because it is a fixed-prize event.
Its implementation cost is that the probability of an exact match can be small when \(n\) is large or the outcome space is rich.
It is a prominent known rule in the comparison, not a new mechanism introduced here.

\subsection{Additional Rules for Computational Comparison}

The next rules broaden the design exercise.
For each rule, the inverse region is fully specified by finite expected-loss inequalities and can therefore be used in the simulation exercise.
What distinguishes these rules from the primary rules is not lack of identification, but the absence of a simple main-text formula for sharp coordinate and mean bounds comparable to the quadratic-distance and discrete-metric cases.
For these rules, the computational representation is the relevant object for the design exercise.

\subsubsection{Manhattan-Distance Frequency Scoring}

The Manhattan-distance rule is
\[
S_M(r,\omega)=a-b\sum_{i=1}^k |r_i-\omega_i|,
\quad b>0.
\]
It rewards total absolute count accuracy rather than squared count accuracy.
Its inferential interpretation is median-oriented rather than mean-oriented.

\paragraph{Optimal reports and belief region.}
For category \(i\), the realized count has binomial marginal distribution
\[
W_i\sim\mathrm{Bin}(n,p_i).
\]
Let
\[
\ell_i(t;p_i)=\E|t-W_i|,
\qquad
F_i(t;n,p_i)=\Pr(W_i\le t).
\]
Then optimal reports solve
\[
R_M(p)
=
\argmin_{r\in\OmegaN}
\sum_{i=1}^k \ell_i(r_i;p_i).
\]
The simplex constraint \(\sum_i r_i=n\) prevents the researcher from treating each coordinate independently.
The relevant optimality condition is that no one-count transfer lowers expected Manhattan loss.
Equivalently,
\[
P_M(r)
=
\left\{
p\in\DeltaK:
F_i(r_i;n,p_i)\ge F_j(r_j-1;n,p_j)
\quad
\forall i,j\text{ with }r_i<n,\ r_j>0
\right\}.
\]
Weak inequalities include ties; there is no receiver constraint when \(r_i=n\) and no sender constraint when \(r_j=0\).

\paragraph{Identification bounds.}
For \(k=2\), write \(r=(t,n-t)\).
Then Manhattan loss is \(2|t-W|\), with \(W\sim\mathrm{Bin}(n,p_1)\), so the report identifies the binomial-median interval
\[
P_{M,1}(r)
=
\{p_1\in[0,1]:F(t-1;n,p_1)\le 1/2\le F(t;n,p_1)\}.
\]
This binary case is clean.

For \(k>2\), the CDF inequalities remain explicit.
The useful computational representation is a one-dimensional threshold search.
Specifically, \(p\in P_M(r)\) if and only if there exists \(c\in[0,1]\) such that
\[
F_i(r_i-1;n,p_i)\le c\le F_i(r_i;n,p_i)
\quad\forall i,
\]
using \(F_i(-1;n,p_i)=0\) and \(F_i(n;n,p_i)=1\).
For a fixed \(c\), these inequalities give coordinate boxes through inverse binomial CDFs; intersecting those boxes with \(\sum_i p_i=1\) gives the feasible beliefs at that threshold.
The sharp coordinate and mean bounds are computed by optimizing over feasible thresholds.

\paragraph{Practical interpretation.}
Manhattan distance is useful when a median-type count prediction is substantively appropriate.
It is also useful as a contrast: changing the distance criterion changes the inferred belief region.
For mean-oriented belief elicitation, however, the multi-category Manhattan bounds are less transparent than the quadratic-distance bounds.

\subsubsection{Hamming-Distance Frequency Scoring}

The Hamming-distance rule penalizes the number of categories whose reported count differs from the realized count:
\[
S_H(r,\omega)=a-b\sum_{i=1}^k\1\{r_i\ne\omega_i\},
\quad b>0.
\]
Its expected loss separates into binomial marginal exact-coordinate probabilities:
\[
\E_p\left[\sum_i \1\{r_i\ne\omega_i\}\right]
=
\sum_i \{1-\Pr(\mathrm{Bin}(n,p_i)=r_i)\}.
\]
Thus the forward objective is computationally simple, even though the reported counts still have to sum to \(n\).

For an observed report \(r\), define
\[
L_H(r;p)
=
\sum_i \{1-\Pr(\mathrm{Bin}(n,p_i)=r_i)\}.
\]
The inverse region is exactly
\[
P_H(r)
=
\{p\in\DeltaK:L_H(r;p)\le L_H(s;p)\ \forall s\in\OmegaN\}.
\]
This finite set of expected-loss inequalities fully specifies identification.
Unlike the primary rules, however, the region is not a simple transfer polytope or a closed coordinate interval formula.
The bounds are therefore computed from the exact inverse-region inequalities in the design exercise.

For \(k=2\), Hamming distance coincides with discrete-metric frequency scoring.
Matching one coordinate is equivalent to matching the whole binary count vector, so the report incentives are the same.

\subsubsection{Chebyshev-Distance Frequency Scoring}

The Chebyshev-distance rule penalizes the largest coordinate-level count error:
\[
S_\infty(r,\omega)=a-b\max_i |r_i-\omega_i|,
\quad b>0.
\]
For \(k=2\), with \(r=(t,n-t)\) and \(\omega=(W,n-W)\),
\[
\|r-\omega\|_\infty=|t-W|.
\]
Manhattan loss is \(2|t-W|\) in the same binary environment.
Thus Chebyshev and Manhattan distance induce the same optimal reports when \(k=2\).

For \(k>2\), the maximum couples coordinates inside each realization.
The expected-loss comparison is therefore nonseparable.
The inverse region is exactly
\[
P_\infty(r)
=
\left\{
p\in\DeltaK:
\E_p\|r-\omega\|_\infty
\le
\E_p\|s-\omega\|_\infty
\quad
\forall s\in\OmegaN
\right\}.
\]
This region is finite and computationally checkable by enumerating feasible reports and evaluating expected losses.
As with Hamming distance, the bounds used in the design exercise are computed from the exact inverse-region inequalities rather than from a simple closed main-text formula.
